\documentclass[10pt,a4paper]{amsart}
\usepackage{amsmath,amssymb,amsthm,graphicx}
% For commutative diagrams we use Xy-pic
\usepackage[all]{xy}

% We use 2cell for 2-commutative diagrams.
\xyoption{2cell}
% Theorem environments.
%
\theoremstyle{plain}
\newtheorem{theorem}[subsection]{Theorem}
\newtheorem{proposition}[subsection]{Proposition}
\newtheorem{lemma}[subsection]{Lemma}

\theoremstyle{definition}
\newtheorem{definition}[subsection]{Definition}
\newtheorem{example}[subsection]{Example}
\newtheorem{exercise}[subsection]{Exercise}
\newtheorem{situation}[subsection]{Situation}

\theoremstyle{remark}
\newtheorem{remark}[subsection]{Remark}
\newtheorem{remarks}[subsection]{Remarks}

\numberwithin{equation}{subsection}

% Macros
%
\def\lim{\mathop{\mathrm{lim}}\nolimits}
\def\colim{\mathop{\mathrm{colim}}\nolimits}
\def\Spec{\mathop{\mathrm{Spec}}}
\def\Hom{\mathop{\mathrm{Hom}}\nolimits}
\def\Ext{\mathop{\mathrm{Ext}}\nolimits}
\def\SheafHom{\mathop{\mathcal{H}\!\mathit{om}}\nolimits}
\def\SheafExt{\mathop{\mathcal{E}\!\mathit{xt}}\nolimits}
\def\Sch{\mathit{Sch}}
\def\Mor{\mathop{\mathrm{Mor}}\nolimits}
\def\Ob{\mathop{\mathrm{Ob}}\nolimits}
\def\Sh{\mathop{\mathit{Sh}}\nolimits}
\def\NL{\mathop{N\!L}\nolimits}
\def\CH{\mathop{\mathrm{CH}}\nolimits}
\def\proetale{{pro\text{-}\acute{e}tale}}
\def\etale{{\acute{e}tale}}
\def\QCoh{\mathit{QCoh}}
\def\Ker{\mathop{\mathrm{Ker}}}
\def\Im{\mathop{\mathrm{Im}}}
\def\Coker{\mathop{\mathrm{Coker}}}
\def\Coim{\mathop{\mathrm{Coim}}}

% Boxtimes
%
\DeclareMathSymbol{\boxtimes}{\mathbin}{AMSa}{"02}

%
% Macros for moduli stacks/spaces
%
\def\QCohstack{\mathcal{QC}\!\mathit{oh}}
\def\Cohstack{\mathcal{C}\!\mathit{oh}}
\def\Spacesstack{\mathcal{S}\!\mathit{paces}}
\def\Quotfunctor{\mathrm{Quot}}
\def\Hilbfunctor{\mathrm{Hilb}}
\def\Curvesstack{\mathcal{C}\!\mathit{urves}}
\def\Polarizedstack{\mathcal{P}\!\mathit{olarized}}
\def\Complexesstack{\mathcal{C}\!\mathit{omplexes}}
% \Pic is the operator that assigns to X its picard group, usage \Pic(X)
% \Picardstack_{X/B} denotes the Picard stack of X over B
% \Picardfunctor_{X/B} denotes the Picard functor of X over B
\def\Pic{\mathop{\mathrm{Pic}}\nolimits}
\def\Picardstack{\mathcal{P}\!\mathit{ic}}
\def\Picardfunctor{\mathrm{Pic}}
\def\Deformationcategory{\mathcal{D}\!\mathit{ef}}
\usepackage{cleveref}
\crefname{originallemma}{original lemma}{original lemmas}
\Crefname{originallemma}{Original lemma}{Original lemmas}
\makeatletter
\newlabel{algebra-lemma-Schanuel}{{00O3}{source}{}{sourcepointer-1}{}}
\newlabel{algebra-lemma-Schanuel@cref}{{[originallemma][0][]00O3}{[1][0][]source}{}{}{}}
\makeatother
\crefname{originallemma}{original lemma}{original lemmas}
\Crefname{originallemma}{Original lemma}{Original lemmas}
\makeatletter
\newlabel{flat-lemma-blowup-complex-integral}{{0ESR}{source}{}{sourcepointer-2}{}}
\newlabel{flat-lemma-blowup-complex-integral@cref}{{[originallemma][0][]0ESR}{[1][0][]source}{}{}{}}
\makeatother
\crefname{originallemma}{original lemma}{original lemmas}
\Crefname{originallemma}{Original lemma}{Original lemmas}
\makeatletter
\newlabel{flat-remark-when-you-have-a-complex}{{0F8J}{source}{}{sourcepointer-3}{}}
\newlabel{flat-remark-when-you-have-a-complex@cref}{{[originallemma][0][]0F8J}{[1][0][]source}{}{}{}}
\makeatother
\crefname{originallemma}{original lemma}{original lemmas}
\Crefname{originallemma}{Original lemma}{Original lemmas}
\makeatletter
\newlabel{lemma-additivity-chern-classes}{{02UI}{source}{}{sourcepointer-4}{}}
\newlabel{lemma-additivity-chern-classes@cref}{{[originallemma][0][]02UI}{[1][0][]source}{}{}{}}
\makeatother
\crefname{originallemma}{original lemma}{original lemmas}
\Crefname{originallemma}{Original lemma}{Original lemmas}
\makeatletter
\newlabel{lemma-bivariant-zero}{{02UC}{source}{}{sourcepointer-5}{}}
\newlabel{lemma-bivariant-zero@cref}{{[originallemma][0][]02UC}{[1][0][]source}{}{}{}}
\makeatother
\crefname{originallemma}{original lemma}{original lemmas}
\Crefname{originallemma}{Original lemma}{Original lemmas}
\makeatletter
\newlabel{lemma-chern-character-additive}{{0F9C}{source}{}{sourcepointer-6}{}}
\newlabel{lemma-chern-character-additive@cref}{{[originallemma][0][]0F9C}{[1][0][]source}{}{}{}}
\makeatother
\crefname{originallemma}{original lemma}{original lemmas}
\Crefname{originallemma}{Original lemma}{Original lemmas}
\makeatletter
\newlabel{lemma-disjoint-decomposition-bivariant}{{0FDZ}{source}{}{sourcepointer-7}{}}
\newlabel{lemma-disjoint-decomposition-bivariant@cref}{{[originallemma][0][]0FDZ}{[1][0][]source}{}{}{}}
\makeatother
\crefname{originallemma}{original lemma}{original lemmas}
\Crefname{originallemma}{Original lemma}{Original lemmas}
\makeatletter
\newlabel{remark-extend-to-finite-locally-free}{{0ESW}{source}{}{sourcepointer-8}{}}
\newlabel{remark-extend-to-finite-locally-free@cref}{{[originallemma][0][]0ESW}{[1][0][]source}{}{}{}}
\makeatother
\crefname{originallemma}{original lemma}{original lemmas}
\Crefname{originallemma}{Original lemma}{Original lemmas}
\makeatletter
\newlabel{situation-setup}{{02QL}{source}{}{sourcepointer-9}{}}
\newlabel{situation-setup@cref}{{[originallemma][0][]02QL}{[1][0][]source}{}{}{}}
\makeatother
\title{Independence of the Chern class, Chern character and Chern polynomials of a bounded finite locally free complex representing a perfect object}
\author{The Stacks Project Authors}
\date{stacks-project:chow.tex:local-snapshot-20261012; GNU FDL 1.2 or later}
\begin{document}
\maketitle
\noindent\textit{Editorial note.} All delivered mathematical text is the author's own, copied from the pinned Stacks Project chapter file chow.tex; no statement or proof step was rewritten, reordered or removed.\par\medskip
\noindent\textit{Editorial note.} Preamble provenance: the pinned chapter file opens with \textbackslash{}input\{preamble\} and carries no preamble text of its own, so the project's single Stacks-wide preamble staged at sources/stacks-chow/preamble.tex is excerpted verbatim as two spans, lines 14-18 (Xy-pic setup) and lines 154-217 (the native theorem declarations and mathematical macros). Nothing was added, removed or redefined; the wrapper class and the amsmath/amssymb/amsthm/graphicx base packages are supplied by the assembler.\par\medskip
\noindent\textit{Editorial note.} Same-chapter pointers are located at their \textbackslash{}label lines inside the pinned chapter file; cross-chapter pointers are located in the project's own tag list staged in the snapshot (sources/stacks-chow/tags.txt), at the verbatim lines carrying their canonical tags.\par\medskip
\noindent\textit{Allowed standard prerequisites (editorial).}\begin{itemize}
\item Additivity of the total Chern class and of the Chern character. Same chapter chow.tex, labels lemma-additivity-chern-classes (used at proof lines 9111 and 9167) and lemma-chern-character-additive (used at proof line 9109). Let the product formula for the total Chern class of a complex be reorganized as a product over kernels and images, and the same for the Chern character; the reductions themselves are inside the counted span.
\item Extension of the construction from finite locally free sheaves to complexes. Same chapter chow.tex, label remark-extend-to-finite-locally-free, invoked at proof line 9114. Fixes what c of a locally bounded complex of finite locally free modules is supposed to mean, so that independence of the representing complex is a well posed question.
\item Splitting of the bivariant theory over an open and closed decomposition, and vanishing of a bivariant class. Same chapter chow.tex, labels lemma-disjoint-decomposition-bivariant (proof line 9131) and lemma-bivariant-zero (proof line 9141). Gives A\textasciicircum{}p(X) as the product of the A\textasciicircum{}p(X\_i) so that the components can be glued, and reduces equality of bivariant classes to equality after capping with [Y] for every integral Y.
\item Tor dimension at most one after a proper birational modification, and Schanuel's lemma. More on Flatness, labels flat-lemma-blowup-complex-integral (proof line 9148) and flat-remark-when-you-have-a-complex (proof line 9163); Algebra, label algebra-lemma-Schanuel (proof line 9203). Reduces to complexes of constant rank whose cohomology has tor dimension at most one, in which case kernels and images are finite locally free, and identifies the two kernels in the final four-term diagram chase.
\end{itemize}\medskip
\ifdefined\hypertarget\hypertarget{sourcepointer-1}{}\fi
\noindent\textit{Source reference.} Original lemma 00O3: sources/stacks-chow/tags.txt, source lines 881--881. Navigation pointer for \textbackslash{}ref\{algebra-lemma-Schanuel\} used inside the counted span (line 9203); the argument that consumes it is retained in the span.\par\medskip
\ifdefined\hypertarget\hypertarget{sourcepointer-2}{}\fi
\noindent\textit{Source reference.} Original lemma 0ESR: sources/stacks-chow/tags.txt, source lines 18491--18491. Navigation pointer for \textbackslash{}ref\{flat-lemma-blowup-complex-integral\} used inside the counted span (line 9148); the argument that consumes it is retained in the span.\par\medskip
\ifdefined\hypertarget\hypertarget{sourcepointer-3}{}\fi
\noindent\textit{Source reference.} Original lemma 0F8J: sources/stacks-chow/tags.txt, source lines 19044--19044. Navigation pointer for \textbackslash{}ref\{flat-remark-when-you-have-a-complex\} used inside the counted span (line 9163); the argument that consumes it is retained in the span.\par\medskip
\ifdefined\hypertarget\hypertarget{sourcepointer-4}{}\fi
\noindent\textit{Source reference.} Original lemma 02UI: sources/stacks-chow/pinned-original-source.tex, source lines 7965--7965. Navigation pointer for \textbackslash{}ref\{lemma-additivity-chern-classes\} used inside the counted span (line 9111); the argument that consumes it is retained in the span.\par\medskip
\ifdefined\hypertarget\hypertarget{sourcepointer-5}{}\fi
\noindent\textit{Source reference.} Original lemma 02UC: sources/stacks-chow/pinned-original-source.tex, source lines 6585--6585. Navigation pointer for \textbackslash{}ref\{lemma-bivariant-zero\} used inside the counted span (line 9141); the argument that consumes it is retained in the span.\par\medskip
\ifdefined\hypertarget\hypertarget{sourcepointer-6}{}\fi
\noindent\textit{Source reference.} Original lemma 0F9C: sources/stacks-chow/pinned-original-source.tex, source lines 8958--8958. Navigation pointer for \textbackslash{}ref\{lemma-chern-character-additive\} used inside the counted span (line 9109); the argument that consumes it is retained in the span.\par\medskip
\ifdefined\hypertarget\hypertarget{sourcepointer-7}{}\fi
\noindent\textit{Source reference.} Original lemma 0FDZ: sources/stacks-chow/pinned-original-source.tex, source lines 6623--6623. Navigation pointer for \textbackslash{}ref\{lemma-disjoint-decomposition-bivariant\} used inside the counted span (line 9131); the argument that consumes it is retained in the span.\par\medskip
\ifdefined\hypertarget\hypertarget{sourcepointer-8}{}\fi
\noindent\textit{Source reference.} Original lemma 0ESW: sources/stacks-chow/pinned-original-source.tex, source lines 7632--7632. Navigation pointer for \textbackslash{}ref\{remark-extend-to-finite-locally-free\} used inside the counted span (line 9114); the argument that consumes it is retained in the span.\par\medskip
\ifdefined\hypertarget\hypertarget{sourcepointer-9}{}\fi
\noindent\textit{Source reference.} Original lemma 02QL: sources/stacks-chow/pinned-original-source.tex, source lines 1372--1372. Navigation pointer for \textbackslash{}ref\{situation-setup\} used inside the counted span (line 9092); the argument that consumes it is retained in the span.\par\medskip
\begin{lemma}
\label{lemma-pre-derived-chern-class}
Let $(S, \delta)$ be as in Situation \ref{situation-setup}.
Let $X$ be locally of finite type over $S$. Let $E \in D(\mathcal{O}_X)$
be an object such that there exists a locally bounded complex
$\mathcal{E}^\bullet$ of finite locally free $\mathcal{O}_X$-modules
representing $E$. Then a slight generalization of the above constructions
$$
c(\mathcal{E}^\bullet) \in \prod\nolimits_{p \geq 0} A^p(X),\quad
ch(\mathcal{E}^\bullet) \in
\prod\nolimits_{p \geq 0} A^p(X) \otimes \mathbf{Q},\quad
P_p(\mathcal{E}^\bullet) \in A^p(X)
$$
are independent of the choice of the complex $\mathcal{E}^\bullet$.
\end{lemma}

\begin{proof}
We prove this for the total Chern class; the other two cases follow
by the same arguments using
Lemma \ref{lemma-chern-character-additive}
instead of
Lemma \ref{lemma-additivity-chern-classes}.

\medskip\noindent
As in Remark \ref{remark-extend-to-finite-locally-free} in order
to define the total chern class $c(\mathcal{E}^\bullet)$
we decompose $X$ into open and closed subschemes
$$
X = \coprod\nolimits_{i \in I} X_i
$$
such that the rank $\mathcal{E}^n$ is constant on $X_i$ for
all $n$ and $i$. (Since these ranks are locally constant functions
on $X$ we can do this.) Since $\mathcal{E}^\bullet$ is locally
bounded, we see that only a finite number of the sheaves
$\mathcal{E}^n|_{X_i}$ are nonzero for a fixed $i$. Hence we
can define
$$
c(\mathcal{E}^\bullet|_{X_i}) =
\prod\nolimits_n c(\mathcal{E}^n|_{X_i})^{(-1)^n}
\in \prod\nolimits_{p \geq 0} A^p(X_i)
$$
as above. By Lemma \ref{lemma-disjoint-decomposition-bivariant}
we have $A^p(X) = \prod_i A^p(X_i)$. Hence for each $p \in \mathbf{Z}$
we have a unique element $c_p(\mathcal{E}^\bullet) \in A^p(X)$ restricting
to $c_p(\mathcal{E}^\bullet|_{X_i})$ on $X_i$ for all $i$.

\medskip\noindent
Suppose we have a second locally bounded complex
$\mathcal{F}^\bullet$ of finite locally free $\mathcal{O}_X$-modules
representing $E$.
Let $g : Y \to X$ be a morphism locally of finite type with $Y$ integral.
By Lemma \ref{lemma-bivariant-zero} it suffices to show that
with $c(g^*\mathcal{E}^\bullet) \cap [Y]$ is the same as
$c(g^*\mathcal{F}^\bullet) \cap [Y]$ and it even suffices to prove
this after replacing $Y$ by an integral scheme proper and birational
over $Y$. Then first we conclude that $g^*\mathcal{E}^\bullet$
and $g^*\mathcal{F}^\bullet$ are bounded complexes of finite locally
free $\mathcal{O}_Y$-modules of constant rank. Next, by
More on Flatness, Lemma \ref{flat-lemma-blowup-complex-integral}
we may assume that $H^i(Lg^*E)$ is perfect of tor dimension $\leq 1$
for all $i \in \mathbf{Z}$.
This reduces us to the case discussed in the next paragraph.

\medskip\noindent
Assume $X$ is integral, $\mathcal{E}^\bullet$ and $\mathcal{F}^\bullet$
are bounded complexes of finite locally free modules of constant rank, and
$H^i(E)$ is a perfect $\mathcal{O}_X$-module of tor dimension $\leq 1$
for all $i \in \mathbf{Z}$. We have to
show that $c(\mathcal{E}^\bullet) \cap [X]$ is the same as
$c(\mathcal{F}^\bullet) \cap [X]$. Denote
$d_\mathcal{E}^i : \mathcal{E}^i \to \mathcal{E}^{i + 1}$ and
$d_\mathcal{F}^i : \mathcal{F}^i \to \mathcal{F}^{i + 1}$
the differentials of our complexes. By
More on Flatness, Remark \ref{flat-remark-when-you-have-a-complex}
we know that $\Im(d_\mathcal{E}^i)$, $\Ker(d_\mathcal{E}^i)$,
$\Im(d_\mathcal{F}^i)$, and $\Ker(d_\mathcal{F}^i)$
are finite locally free $\mathcal{O}_X$-modules for all $i$.
By additivity (Lemma \ref{lemma-additivity-chern-classes}) we see that
$$
c(\mathcal{E}^\bullet) = \prod\nolimits_i
c(\Ker(d_\mathcal{E}^i))^{(-1)^i} c(\Im(d_\mathcal{E}^i))^{(-1)^i}
$$
and similarly for $\mathcal{F}^\bullet$. Since we have the
short exact sequences
$$
0 \to \Im(d_\mathcal{E}^i) \to \Ker(d_\mathcal{E}^i) \to H^i(E) \to 0
\quad\text{and}\quad
0 \to \Im(d_\mathcal{F}^i) \to \Ker(d_\mathcal{F}^i) \to H^i(E) \to 0
$$
we reduce to the problem stated and solved in the next paragraph.

\medskip\noindent
Assume $X$ is integral and we have two short exact sequences
$$
0 \to \mathcal{E}' \to \mathcal{E} \to \mathcal{Q} \to 0
\quad\text{and}\quad
0 \to \mathcal{F}' \to \mathcal{F} \to \mathcal{Q} \to 0
$$
with $\mathcal{E}$, $\mathcal{E}'$, $\mathcal{F}$, $\mathcal{F}'$
finite locally free. Problem: show that
$c(\mathcal{E})c(\mathcal{E}')^{-1} \cap [X] =
c(\mathcal{F})c(\mathcal{F}')^{-1} \cap [X]$.
To do this, consider the short exact sequence
$$
0 \to \mathcal{G} \to \mathcal{E} \oplus \mathcal{F} \to \mathcal{Q} \to 0
$$
defining $\mathcal{G}$. Since $\mathcal{Q}$ has tor dimension $\leq 1$
we see that $\mathcal{G}$ is finite locally free. A diagram chase
shows that the kernel of the surjection $\mathcal{G} \to \mathcal{F}$
maps isomorphically to $\mathcal{E}'$ in $\mathcal{E}$ and
the kernel of the surjection $\mathcal{G} \to \mathcal{E}$ maps
isomorphically to $\mathcal{F}'$ in $\mathcal{F}$. (Working affine
locally this follows from or is equivalent to Schanuel's lemma, see
Algebra, Lemma \ref{algebra-lemma-Schanuel}.)
We conclude that
$$
c(\mathcal{E})c(\mathcal{F}') = c(\mathcal{G}) =
c(\mathcal{F})c(\mathcal{E}')
$$
as desired.
\end{proof}

\end{document}
